% V32 standalone theoretical method note, no real target outcomes included yet.
% Not a BSPC accepted method and not independently registered.
\section{Finite-source contrast uncertainty as a regularisation operator}
\label{sec:v32-theory}
\paragraph{Hypothesis and prior art.}
Anisotropic regularisation, discriminant analysis, random-effects
variance components and source-domain generalisation have substantial
prior literature. Our specific falsifiable construction combines
source-participant contrast variance \emph{deconvolved for finite
within-person trial noise} with the \emph{same} fixed
EEG pooled-LDA classifier, and tests whether the resulting matrix
operator predicts performance shifts under source-only selection.
We do not claim this combination is independently novel until a
thorough literature review and new EEG benefit are established.
A complementary research question is whether the observed
between-person dispersion represents physical discriminant-direction
shift or is predominantly the variance of noisy estimated contrasts.

\paragraph{Source-only estimation.}
Let each independent source recording group $g=1,\dots,G$
contain class-specific trials with means
$\widehat\mu_{g0},\widehat\mu_{g1}$ and class covariance
$\widehat\Sigma_{gc}$. Define the estimated class contrast
$\widehat d_g=\widehat\mu_{g1}-\widehat\mu_{g0}$
and sampling-noise variance estimate
\[
  \widehat V_g=
  {\widehat\Sigma_{g0}\over n_{g0}}+
  {\widehat\Sigma_{g1}\over n_{g1}},\qquad
  \widehat B_{\rm obs}=\operatorname{Cov}_g(\widehat d_g).
\]
The proposed heterogeneity matrix is
\[
  \widehat B_{\rm debias}=
  \Pi_{\mathbb S_+}\left(
       \widehat B_{\rm obs}-{1\over G}\sum_g\widehat V_g
                       \right)
\]
with $\Pi_{\mathbb S_+}$ denoting the standard eigenvalue-wise
positive semidefinite projection. The source-only operator is
\[
 \boxed{\quad
 w_{\lambda,\alpha}=
  \left[(1-\lambda)\widehat S_{\rm pooled}
        +\lambda\nu I+\alpha\widehat B_{\rm debias}
  \right]^{-1}\widehat d_{\rm pooled},
 \quad \nu={\operatorname{tr}(\widehat S_{\rm pooled})\over p}.
 \quad}
\]
The original EEG pipeline is recovered EXACTLY when $\alpha=0$;
with $\widehat B_{\rm debias}=0$ it is identical for all $\alpha$.
Conversely, a non-isotropic matrix may produce an oblique
new direction even at $\lambda=1$ when $\alpha>0$,
avoiding the V31 fixed-$\lambda=1$ endpoint degeneracy.
This is a testable distinction, not a promise of an improvement.

\paragraph{Proposition 1: sampling-noise decomposition.}
Assume independent source-group sampling, unbiased class-mean
estimates, finite second moments and uncorrelated zero-mean
within-group estimation errors
$\widehat d_g=d_g+\varepsilon_g$ with
$\operatorname{Cov}(\varepsilon_g)=V_g$.
Conditionally on the fixed group contrasts $d_1,\ldots,d_G$,
\[
  \mathbb E[\widehat B_{\rm obs}]
  = B_{\rm actual} + {1\over G}\sum_g V_g,
  \qquad
  B_{\rm actual}={1\over G-1}\sum_g(d_g-\bar d)
                                   (d_g-\bar d)^\top .
\]
Thus subtracting the average \emph{unprojected} mean-noise covariance
is unbiased for $B_{\rm actual}$, provided the sample class
covariance estimates are unbiased. Note that the PSD projection is
nonlinear and introduces upward bias and potential direction changes.
No finite-source posterior or target coverage follows from this
identity, especially when $G$ is small compared with feature rank.

\paragraph{Proposition 2: motivation under a random-effect Gaussian target.}
Under \emph{additional assumptions} that target class-conditional
features are Gaussian with shared known covariance $\Sigma\succ0$
and random target class contrast $d_t\sim
\mathcal N(\bar d,B)$ independent of the test noise,
the unconditional expected target AUC of a fixed nonzero direction $w$
is exactly
\[
 \mathbb E_{d_t}\bigl[\operatorname{AUC}_t(w)\bigr]
  =\Phi\!\left(
     {w^\top\bar d\over\sqrt{2w^\top\Sigma w+w^\top Bw}}\right).
\]
To see this, the target positive-minus-negative score is normal
with conditional mean $w^\top d_t$ and conditional variance
$2w^\top\Sigma w$; marginalizing the Gaussian $d_t$ adds
$w^\top Bw$ to the variance. Maximizing its normalized mean
gives the Fisher solution
$w_{\rm rand}\propto(2\Sigma+B)^{-1}\bar d$,
equivalently
$(\Sigma+\tfrac12 B)^{-1}\bar d$.
This is a classical Gaussian random-effects calculation, NOT a
generalization theorem for real EEG. Our $\lambda=0,\alpha=0.5$
special case resembles this reference only if
$\widehat S_{\rm pooled}\simeq\Sigma$,
$\widehat d_{\rm pooled}\simeq\bar d$, and
$\widehat B_{\rm debias}\simeq B$.
If target directions are not drawn from this family, it can be worse
than scalar shrinkage.

\paragraph{Source-only selection and falsification.}
The frozen code evaluates every feasible original 27-point scalar
$\lambda$ at $\alpha\in\{0,0.5,2\}$ on original round-robin
grouped source folds, refitting tangent reference and all
variance components within source fold training. It selects by
group macro-AUC minus $0.25$ times the source-group AUC standard
deviation. The unpenalized grouped source macro-AUC at
$\alpha=0$ is reproduced bitwise to ordinary floating tolerance.
We additionally report the risk-only scalar ablation and
the best source mean-AUC matrix ablation. Source group CV
is a heuristic: its standard deviation is not an uncertainty
interval on unseen persons.

\paragraph{Novelty disqualification and first external study.}
The method fails its \emph{structural} novelty gate if it selects
$\alpha=0$, if $\widehat B_{\rm debias}=0$, or if the selected
direction remains parallel to the corresponding scalar-LDA
direction. Its \emph{empirical} method-promotion gate is
separately frozen in the Dreyer2023A protocol and requires
a material paired improvement, valid sign-flip result and
no increase in adverse held subjects.
We retain the source-only controls, Riemannian MDM and
tangent logistic regression under the same electrodes,
plus full source choice SHA receipts before any held
Dreyer2023A target access. A successful test on one
collection would require independent-acquisition replication
and a prior-art novelty review before any L8/L9 claim.
